\documentclass[12pt,a4paper]{article}
        \makeatletter
        \def\input@path{{../}}
        \makeatother
        \usepackage{almeria-fichas}

        \FichaMeta{Sesion 19}{Ensayo y mejora}{Bloque 4 · Producto final}{Ensayar la exposicion, aplicar la rubrica y corregir los puntos debiles de la guia.}{Procedimientos}{Analizar, Evaluar, Crear}

        \begin{document}

        \FichaCabecera

        \begin{materialesbox}
        \begin{itemize}
    \item Ficha impresa
    \item Lapiz
    \item Guia casi final
    \item Rubrica de autoevaluacion y coevaluacion
\end{itemize}
        \end{materialesbox}

        \begin{pasosbox}
        \begin{enumerate}
    \item Haz un ensayo corto respetando el tiempo.
    \item Escucha el feedback con atencion y sin interrumpir.
    \item Marca tres mejoras prioritarias.
    \item Aplica los cambios antes de la entrega final.
\end{enumerate}
        \end{pasosbox}

        \begin{recordatoriobox}
        \begin{itemize}
    \item Una buena revision es concreta y respetuosa.
    \item Las mejoras deben basarse en criterios, no en gustos sin justificar.
    \item Corregir tambien es parte del trabajo matematico.
\end{itemize}
        \end{recordatoriobox}

        \begin{apoyobox}
        \begin{itemize}
    \item Puedes usar la estructura: estrella, estrella, deseo.
    \item Escribe primero el problema y despues la solucion.
    \item Si una mejora afecta a una formula o unidad, compruebala dos veces.
\end{itemize}
        \end{apoyobox}

        \BloomSection{analizar}

\BloomExercise{analizar}{1}{Partes y relaciones}{Analiza un ejemplo relacionado con \emph{Ensayo y mejora}. Separa sus partes, indica cuales son relevantes y explica como se relacionan entre si.}{8}

\BloomExercise{analizar}{2}{Detecta errores o diferencias}{Compara dos respuestas, dos representaciones o dos procedimientos relacionados con esta sesion. Analiza sus diferencias y explica que errores o mejoras detectas.}{8}

\BloomSection{evaluar}

\BloomExercise{evaluar}{1}{Puntos debiles}{Evalua tu propia guia con una rubrica sencilla y detecta tres puntos debiles.}{8}

\BloomExercise{evaluar}{2}{Feedback a otro equipo}{Evalua una guia ajena y escribe dos puntos fuertes y dos mejoras concretas.}{8}

\BloomSection{crear}

\BloomExercise{crear}{1}{Plan de mejora}{Crea un plan de mejora con tres cambios: que cambiaras, por que y quien lo hara.}{8}

\BloomExercise{crear}{2}{Version corregida}{Redacta una version corregida de una parte de la guia que antes estaba incompleta o poco clara.}{8}

        \begin{evidenciabox}
        \begin{itemize}
    \item Autoevaluar la guia con criterio.
    \item Aplicar correcciones concretas y justificadas.
    \item Mejorar la exposicion oral y el producto escrito.
\end{itemize}
        \end{evidenciabox}

        \end{document}
